\documentclass[twocolumn]{article}
\usepackage{graphicx}
\usepackage{booktabs}
\usepackage{amsmath}
\usepackage{geometry}
\usepackage{hyperref}
\usepackage{xurl}
\usepackage{array}
\usepackage{float}
\usepackage{multirow}
\geometry{a4paper, margin=0.75in}

\hypersetup{
    colorlinks=true,
    linkcolor=blue,
    filecolor=magenta,
    urlcolor=blue,
    citecolor=blue,
}

\title{\textbf{Nagar Rakshak: Data-Centric Pothole Detection and Severity Pipeline with Ultralytics YOLO Customization}}
\author{NGR-YOLOv8 Project Team}
\date{March 2026}

\begin{document}
\maketitle

\begin{abstract}
\noindent
This paper documents the complete research and engineering progress of the NGR-YOLOv8 project up to the current state of development. The system is designed as a two-stage pipeline: M1 for road-defect detection and M2 for severity classification. Beyond standard YOLOv8 training, this work includes substantial data-centric redesign, runtime integration updates, and direct source-code customization of the local Ultralytics library. We report all currently available quantitative results from saved run artifacts and project logs, including multi-class and single-class detector phases, checkpoint recovery experiments, and M2 training status. The strongest recorded detection result in project logs is from the refined single-class setting (precision 94.7\%, recall 91.6\%, mAP@0.5 95.7\%, mAP@0.5:0.95 81.5\%). The best result in currently available run CSV artifacts is from \texttt{broken\_street\_v8x4\_recover3} (precision 78.3\%, recall 56.3\%, mAP@0.5 64.4\%, mAP@0.5:0.95 35.7\%). We further provide a full change-log of modifications made on top of base Ultralytics, including Lion optimizer integration, Coordinate Attention modules, Mish activation support, and MPDIoU metric support.
\end{abstract}

\section{Introduction}
Road-surface defect analytics is a high-impact application of real-time computer vision, where operational value depends not only on raw recall but also on low false-positive rates and stable deployment behavior. This project began with a multi-class YOLOv8 detector and evolved into a broader research and engineering effort involving:
\begin{itemize}
    \item dataset consolidation from heterogeneous sources,
    \item data-centric false-positive reduction using negative samples,
    \item M1+M2 pipeline design for detection plus severity,
    \item post-training feedback loops,
    \item and direct modifications to Ultralytics internals.
\end{itemize}

The objective of this paper is archival: to record exactly what has been done, what results were obtained, what source-level changes were introduced into the local Ultralytics fork, and what components remain incomplete.

\section{System Overview}
\subsection{Current Runtime Pipeline}
The deployed FastAPI service (\texttt{main.py}) implements the following workflow:
\begin{enumerate}
    \item receive uploaded image,
    \item run M1 detector and produce bounding boxes,
    \item run M2 severity classification on selected crops when an M2 checkpoint is available,
    \item return two overlays (M1 and M2 visualizations), per-detection metadata, and counters.
\end{enumerate}

The runtime model-selection logic prioritizes recent trained checkpoints under \texttt{runs/detect/**/weights/best.pt} and falls back to default pretrained weights only if no project checkpoints exist.

\subsection{M1 and M2 Roles}
\begin{itemize}
    \item \textbf{M1 (Detection):} YOLO-based bounding-box detection of road defects (historically either two classes or single class depending on dataset version).
    \item \textbf{M2 (Severity):} YOLO-based severity classification of detected crops into levels (minor, medium, major), integrated as a second-stage module when an M2 checkpoint is available.
\end{itemize}

\section{Datasets and Data Engineering}
\subsection{Base Dataset Construction}
The initial dataset builder (\texttt{prepare\_dataset.py}) integrates:
\begin{itemize}
    \item archive2 annotations and images,
    \item DatasetNinja pothole data,
    \item additional road-scene imagery from Roboflow exports.
\end{itemize}
All inputs are mapped into YOLO format with class mapping:
\begin{itemize}
    \item class 0: broken street,
    \item class 1: legacy auxiliary class used in earlier multi-class experiments.
\end{itemize}

\subsection{Dataset Split Ratios (Train/Validation/Test)}
The project uses YOLO-style \textbf{train/validation} organization for detector development. In multiple scripts, the default random split is:
\[
\text{Train} : \text{Validation} = 80:20
\]

Concrete split behavior is as follows:
\begin{itemize}
    \item \textbf{archive2 in \texttt{prepare\_dataset.py}:} if \texttt{archive/splits.json} contains entries, those predefined train/test lists are respected; otherwise a random 80/20 assignment is applied.
    \item \textbf{DatasetNinja in \texttt{prepare\_dataset.py}:} random 80/20 split into train and val.
    \item \textbf{Roboflow street-light data in \texttt{prepare\_dataset.py}:} source \texttt{train} is mapped to project train, while source \texttt{valid} and \texttt{test} are both mapped to project val.
    \item \textbf{Refined background addition in \texttt{prepare\_background\_dataset.py}:} background images are shuffled and split 80/20 (reported as 520 train and 131 val, total 651 negatives).
    \item \textbf{Severity dataset in \texttt{prepare\_severity\_dataset.py}:} cropped severity samples use random 80/20 split with \texttt{TRAIN\_SPLIT=0.8}.
\end{itemize}

Important evaluation note: for M1 detector experiments, there is \textbf{no separate held-out test directory in the core YAML files}; model-selection metrics are primarily validation-based (Precision, Recall, mAP@0.5, mAP@0.5:0.95).

\subsection{Refined Data Strategy for False-Positive Reduction}
The refined dataset pipeline (\texttt{prepare\_background\_dataset.py}) applies two major changes:
\begin{enumerate}
    \item remove class 1 labels to form single-class detection,
    \item inject negative/background road images (clean/smooth road) with no labels.
\end{enumerate}

Project logs report \textbf{651 negative images} added (520 train, 131 val).

\subsection{Feedback-Driven Fine-Tuning Data}
The script \texttt{fine\_tune\_feedback.py} builds \texttt{combined\_dataset} by:
\begin{itemize}
    \item copying refined base training/validation splits,
    \item appending false-positive feedback images,
    \item oversampling feedback items by a factor of 50.
\end{itemize}
This enables hard-negative focused adaptation without relabeling the full corpus.

\section{Training Configuration History}
\subsection{Detection Training Phases}
Table \ref{tab:training_phases} summarizes the key M1 training phases.

\begin{table}[H]
\centering
\caption{M1 Training Phases and Configuration Evolution}
\label{tab:training_phases}
\resizebox{\columnwidth}{!}{%
\begin{tabular}{@{}lllll@{}}
\toprule
\textbf{Run/Phase} & \textbf{Classes} & \textbf{Backbone} & \textbf{Notable Settings} & \textbf{Status} \\ \midrule
\texttt{broken\_street\_v8x4} & 2 & YOLOv8x & AdamW, imgsz 640, batch 4 & Interrupted/unstable tail \\
\texttt{broken\_street\_v8x4\_recover3} & 2 & YOLOv8x (from best.pt) & imgsz 512, batch 2, workers 0, mosaic/mixup/copy\_paste disabled & Completed (12 epochs) \\
\texttt{refined\_pothole\_v8m} (reported) & 1 & YOLOv8m & AdamW, imgsz 640, batch 8, refined negatives & Completed in report \\
\texttt{pothole\_finetuned\_feedback} & 1 & warm-start from best & feedback oversampling, lr0 5e-4, 30 epochs & Scripted/iterative \\
\bottomrule
\end{tabular}%
}
\end{table}

\subsection{Severity Training Variants}
M2 experimentation is maintained as a YOLO-only severity track:
\begin{itemize}
    \item YOLO classification experiments on cropped pothole regions,
    \item severity-label data preparation from M1 detections for minor/medium/major classes,
    \item application-layer M1+M2 integration for conditional severity inference.
\end{itemize}

\subsection{Reproducible Training Setup}
The practical end-to-end setup used in this project is:
\begin{enumerate}
    \item install Python dependencies from \texttt{requirements.txt},
    \item download/acquire raw datasets (Roboflow exports, archive2, DatasetNinja),
    \item generate YOLO detection dataset via \texttt{prepare\_dataset.py},
    \item optionally generate refined single-class dataset with negatives via \texttt{prepare\_background\_dataset.py},
    \item train detector using one of \texttt{train.py}, \texttt{train\_refined.py}, or \texttt{resume\_recover.py},
    \item prepare severity crops via \texttt{prepare\_severity\_dataset.py}, then train M2 using \texttt{train\_severity.py} or \texttt{train\_m2\_custom.py}.
\end{enumerate}

Environment and stability adjustments recorded in project reports include:
\begin{itemize}
    \item PyTorch 2.6.0 with CUDA 12.4,
    \item reduced dataloader and augmentation pressure during recovery runs (workers=0, mosaic/mixup/copy\_paste disabled),
    \item checkpoint recovery by restarting from \texttt{best.pt} when \texttt{last.pt} was corrupted.
\end{itemize}

\subsection{What Changed Across Training Iterations}
Key model-development changes across the project timeline are summarized below:
\begin{itemize}
    \item \textbf{Taxonomy change:} moved from a two-class detector (broken street + street light) to a focused single-class detector for road defects in refined training.
    \item \textbf{Data change:} added large-scale good-road negative samples to suppress false positives.
    \item \textbf{Recovery change:} after resume instability, shifted to recovery training from stable \texttt{best.pt} with memory-safe hyperparameters.
    \item \textbf{Feedback loop change:} introduced hard-negative fine-tuning with 50$\times$ oversampling of user feedback samples.
    \item \textbf{System change:} integrated dynamic latest-checkpoint loading and optional M2 severity inference in the API pipeline.
\end{itemize}

\section{Results}
\subsection{Results from Available Run CSV Artifacts}
Table \ref{tab:csv_results} reports best values directly extracted from saved CSV logs in \texttt{runs/detect/runs/detect}.

\begin{table}[H]
\centering
\caption{Verified Metrics from Available CSV Runs}
\label{tab:csv_results}
\resizebox{\columnwidth}{!}{%
\begin{tabular}{@{}lcccc@{}}
\toprule
\textbf{Run} & \textbf{Precision} & \textbf{Recall} & \textbf{mAP@0.5} & \textbf{mAP@0.5:0.95} \\ \midrule
\texttt{broken\_street\_v8x4} (best observed) & 74.8\% & 59.2\% & 65.0\% & 34.5\% \\
\texttt{broken\_street\_v8x4\_recover3} (best observed) & 78.3\% & 56.3\% & 64.4\% & 35.7\% \\
\texttt{broken\_street\_v8x4\_recover3} (epoch 12 final) & 75.8\% & 56.8\% & 63.9\% & 35.3\% \\
\bottomrule
\end{tabular}
}
\end{table}

\subsection{Reported Single-Class Metrics}
Project status logs and prior progress reports record the following result for \texttt{refined\_pothole\_v8m}:
\begin{itemize}
    \item precision 94.7\%,
    \item recall 91.6\%,
    \item mAP@0.5 95.7\%,
    \item mAP@0.5:0.95 81.5\%.
\end{itemize}

These metrics are retained as \textbf{reported historical results}. The corresponding run directory is not present in currently available run folders, so this paper marks them as report-derived rather than artifact-reparsed.

\subsection{Qualitative Findings}
Across logs and integration tests:
\begin{itemize}
    \item API endpoints (/ and /predict) were confirmed operational,
    \item detector inference returns class labels and confidences correctly,
    \item sample uploads produced positive detections after model-selection fixes,
    \item M2 inference is conditional and depends on severity checkpoint availability.
\end{itemize}

\section{Ultralytics Library Modifications}
This project includes direct source-code edits to the local Ultralytics fork. Table \ref{tab:ultra_changes} summarizes all confirmed modifications against the vendored repository status.

\begin{table*}[t]
\centering
\caption{Confirmed Source-Level Changes Made to Local Ultralytics Library}
\label{tab:ultra_changes}
\footnotesize
\setlength{\tabcolsep}{4pt}
\begin{tabular}{@{}>{\raggedright\arraybackslash}p{0.28\textwidth}>{\raggedright\arraybackslash}p{0.68\textwidth}@{}}
\toprule
\textbf{File} & \textbf{Change Summary} \\ \midrule
\path{ultralytics/engine/trainer.py} & Added \texttt{Lion} to supported optimizer set; added optimizer-argument branch for Lion (lower LR scaling); integrated custom Lion class loading from \texttt{ultralytics.utils.optimizers}. \\
\path{ultralytics/utils/optimizers.py} & New file implementing Lion optimizer (EvoLved Sign Momentum) with momentum-sign update rule and optional weight decay. \\
\path{ultralytics/nn/modules/block.py} & Added \texttt{CoordAtt} module (Coordinate Attention) and \texttt{C2fCoordAtt} block for enhanced spatial attention in backbone stages. \\
\path{ultralytics/nn/modules/__init__.py} & Exported new modules \texttt{CoordAtt} and \texttt{C2fCoordAtt}. \\
\path{ultralytics/nn/tasks.py} & Registered \texttt{CoordAtt} and \texttt{C2fCoordAtt} in model parser sets so YAML model configs can instantiate them. \\
\path{ultralytics/nn/modules/conv.py} & Added \texttt{Mish} activation class and enabled \texttt{act="Mish"} branch in convolution module constructor. \\
\path{ultralytics/utils/metrics.py} & Extend